% LSPV-4 REPORT::SOURCELDEGATE2::2026-10-14T11:36:00+09:00
\documentclass[11pt]{article}
\usepackage{amsmath,amssymb,amsthm,hyperref}
\newtheorem{theorem}{Theorem}
\title{SOURCELDEGATE2: Source-Extension and Implicit-Tree Time Gates}
\author{LSPV-4, Nonrelativizing Simulation}
\date{14 October 2026}
\begin{document}\maketitle
\noindent\textbf{P versus NP remains OPEN.} Results are exact only in
their declared models; independent review pending.
\section{IDEAS PROPOSED BEFORE LITERATURE}
Candidate A: unrestricted fast explicit-source multilinear extension
evaluation might already imply fast \#SAT. Test repeated inputs.
Candidate B: read-once structure after conditioning on all repeated
variables might give a polynomial extension evaluator when that
repetition set is logarithmic. Test finite fields exhaustively.
Candidate C: the proposed square-root-space simulator uses the
original child time exponent, not the original child space exponent.
Audit the same-machine recurrence.
Candidate D: polynomial time on an explicit Tree Evaluation instance
might conceal exponential expansion of an implicit dependency DAG.
Test ordered-port duplication and the 2026 catalytic results.
\section{LITERATURE CHECK AND REVISIONS}
Williams (STOC 2025) proves a multitape
$\mathrm{TIME}[t]\subseteq\mathrm{SPACE}[O(\sqrt{t\log t})]$
simulation without the time guarantee needed here. His Remark 1.6
excludes arbitrary random access; Theorem 5.1 retains a query-length
restricted oracle version. Section 5 already warns about fast
canonical multilinear extensions. Thus A is not claimed novel.
Cook--Mertz (STOC 2024) supplies the original tree algorithm.
Asadi--Cleve (arXiv:2604.02606, 2026) give
$n^{1/\varepsilon+O(1)}$-time explicit Tree Evaluation with almost
logarithmic total space; Henzinger--Pyne--Ragavan
(arXiv:2602.14320) give another catalytic tradeoff.
Neither theorem supplies a near-linear-time algorithm for Williams'
implicitly unfolded tree as a function of original $t$.
The FLMV (JACM 2005) forbidden region is
$\beta(c)=((c+2)/c^2-c)/2$ on their random-access model.
\section{CONSTRUCTION OR PROOF ATTEMPT}
Let $C$ be an explicit Boolean formula tree of size $s$ with input
variables $x_1,\ldots,x_n$. Let $R(C)$ contain variable names
appearing at least twice among its leaves; set $r=|R(C)|$.
For any finite field, define its multilinear extension
\[
 M_C(a)=\sum_{u\in\{0,1\}^n}C(u)
 \prod_{i:u_i=1}a_i\prod_{i:u_i=0}(1-a_i).
\]
\begin{theorem}[Explicit-source repetition gate]
$M_C(a)$ can be evaluated using $O(2^r(s+r))$ field operations
and polynomial workspace. In particular $r=O(\log s)$ gives
polynomial bit time if field operations have polynomial bit cost.
\end{theorem}
\begin{proof}
Enumerate the $2^r$ assignments $z$ to $R(C)$. Substitution makes
the remaining formula syntactically read-once. At disjoint-support
gates, compute its polynomial expectation by
$\neg P=1-P$, $P\wedge Q=PQ$, $P\vee Q=P+Q-PQ$,
$P\oplus Q=P+Q-2PQ$. These identities hold because the child
supports are disjoint. Multiply by
$\prod_{i\in R(C)}(z_i?a_i:1-a_i)$ and sum over $z$.
This is exactly the defining interpolation sum, with the
variables partitioned into repeated and nonrepeated sets.
Each residual requires one tree traversal, and sequential
evaluation reuses a polynomial-size stack. \end{proof}
This is a formula-tree theorem, not an uncharged DAG expansion.
The rule $P\wedge Q=PQ$ fails with repeated support: for
$x\wedge x$ at $a=2$, the naive result is $4$ but
$M_{x\wedge x}(2)=2$ in a field of odd characteristic.
\begin{theorem}[Exact counting gate]
A uniform deterministic polynomial-bit-time evaluator of
$M_C(1/2,\ldots,1/2)$ for arbitrary Boolean formulas $C$,
over arbitrary supplied odd prime fields, computes \#SAT in
polynomial time and hence implies $P=NP$.
\end{theorem}
\begin{proof}
In field $\mathbb F_p$,
$M_C(1/2,\ldots,1/2)=\#\mathrm{SAT}(C)\,2^{-n}$.
Evaluate over the first $n+1$ odd primes, multiply by $2^n$,
and reconstruct the integer in $[0,2^n]$ by Chinese remaindering.
The primes have polynomial magnitude and can be enumerated
in polynomial time. Their product exceeds $2^n$.
All input and output bit lengths are polynomial. \end{proof}
This does not prohibit Cook--Mertz's nonmultilinear low-degree
extensions, whose representation is distinct.
\begin{theorem}[Implicit-tree transfer barrier]
An ordered-input binary DAG with $h+1$ labelled nodes and two
ports from $v_{j-1}$ to $v_j$ has an unfolded tree with exactly
$2^{h+1}-1$ node occurrences. A generic algorithm polynomial
in explicit tree size therefore yields no polynomial-in-$h$
bound for succinctly represented trees without a separate theorem.
\end{theorem}
\begin{proof}
At depth $j$ there are $2^j$ occurrences of the source node.
Sum the geometric series. The DAG is itself easy to evaluate;
the result is a model-transfer warning, not a hardness claim.
\end{proof}
\paragraph{Joint-resource correction.}
If a hypothetical native model-matched simulator of time $t$ costs
$t^{q+o(1)}$ time and $t^{1/2+o(1)}$ space, then applying it
to $N^{3/5}$ materialized children of length $N^{1/2}$, with
$N^{11/10}$ local work and $N^{1/2}$ live parent space, gives
\[
 \tau'=\max\{11/10,3/5+q\tau/2\},\qquad
 \sigma'=\max\{1/2,\tau/4\}.
\]
The old $\sigma/2$ child-space recurrence does not describe this
simulated realization. For $q<2$ the limiting time exponent is
$c_*=6/[5(2-q)]$ in the relevant range. The limiting space
exponent is $1/2$ if $c_*<2$. The FLMV sufficient frontier yields
$q<1.0046197508498036\ldots$ as a conditional threshold.
No such simulator or all-input contractor is known.
\section{TESTS AND OUTCOME}
One V8 batch with deterministic seed 20261014 compared formula
extension arithmetic against independent full Boolean truth tables:
240 read-once, 336 conditioned-overlap, 128 modular-CRT counting,
one deliberate repeated-variable failure, 24 tree-unfolding, and two
frontier sign checks, total \textbf{731 assertions, zero failures}.
At $q=1.0045$, $\beta(c_*)=0.50028797>1/2$; at $q=1.005$,
$\beta(c_*)=0.49908562<1/2$. A height-24 ordered-port tree
has 33,554,431 occurrences but only 25 DAG nodes.
No SAT lower bound, asymptotic runtime experiment or full
Cook--Mertz implementation is claimed.
\section{Barrier audit and handoff}
These results provide source-sensitive fast and hard
\emph{canonical-extension} cases, not a genuinely new
nonrelativizing SAT separation. The ordinary simulation of Williams
and its restricted oracle extension must remain separate.
COMPLETIONCUT1 and LIVEPREFIX1 preserve the distinction between
pre-release prediction and post-release computation. No hash,
cryptographic assumption, expanded proof-system lower bound,
or Natural-Proofs claim is used.
The key next target is fast evaluation of the specific
nonmultilinear block-transition extensions in the Williams
Tree Evaluation reduction, with native simultaneous time and
space bounds and explicit random-access conversion.
The root \texttt{main.tex} remains truncated and without balanced
LSPV-4 ownership markers. We have not modified it. The standalone
report and code/JSON publication require exact remote readback;
failed writes remain PENDING\_SYNC.
\begin{thebibliography}{6}
\bibitem{W} R. Williams, Simulating Time With Square-Root Space,
STOC 2025.
\bibitem{CM} J. Cook and I. Mertz, Tree Evaluation Is in Space
$O(\log n\log\log n)$, STOC 2024.
\bibitem{AC} V. Asadi and R. Cleve, arXiv:2604.02606, 2026.
\bibitem{HPR} A. Henzinger, E. Pyne and S. Ragavan,
arXiv:2602.14320, 2026.
\bibitem{V} L. Valiant, SIAM J. Comput. 8, 410--421, 1979.
\bibitem{F} L. Fortnow, R. Lipton, D. van Melkebeek,
A. Viglas, J. ACM 52, 835--865, 2005.
\end{thebibliography}
\end{document}
% END_OF_VERIFIED_REPORT::LSPV-4::SOURCELDEGATE2
